\chapter{系统设计方法论}

系统设计轮没有标准答案，但有标准\emph{打法}。这一章给你一套可以套用在任何题目上的五步框架，以及背后的估算与一致性基础。后面三章（订房、搜索、消息/支付）都会严格按这套框架来组织。

\section{五步答题框架}

\begin{enumerate}
  \item \textbf{澄清需求（5 分钟）}：别急着画图。先问清楚——
  \begin{itemize}
    \item \emph{功能性}：核心用例是什么？哪些先做、哪些可以砍？
    \item \emph{非功能性}：日活/QPS 量级？读写比？延迟要求（p95/p99）？一致性要求（强一致还是最终一致即可）？
  \end{itemize}
  \item \textbf{容量估算（3 分钟）}：用数量级估算 QPS、存储、带宽。目的是为后面的分片、缓存决策提供依据，不是为了精确。
  \item \textbf{高层架构（8 分钟）}：画出主要组件——客户端、API 网关、服务、缓存、数据库、消息队列、搜索引擎。强调\textbf{解耦}与\textbf{可扩展}。
  \item \textbf{数据模型 + API（8 分钟）}：给出关键实体的表结构和核心 API 契约。
  \item \textbf{深挖 + 瓶颈/容错（15 分钟）}：挑 1--2 个核心难点深入（通常是并发一致性或热点），再谈扩展、缓存、限流、容灾。
\end{enumerate}

\begin{idea}[“先做简单版本，再演进”]
面试官最想看到的叙事是：\textbf{Start simple, then evolve}。先给一个能跑通的单体/单库方案，再主动指出它的瓶颈（“当 QPS 到 X，这里会成为热点”），然后引入缓存/分片/队列来解决。这种“演进式”叙事远胜于一上来就堆满 Kafka、Redis、分库分表却说不清为什么。复杂度要由\emph{问题}逼出来，而不是为了炫技。
\end{idea}

\section{容量估算速查}

\begin{keypoint}[常用数量级]
$1$ 天 $\approx 86{,}400$ 秒 $\approx 10^5$ 秒。所以“日 $1$ 亿次请求”$\approx 1000\ \text{QPS}$（均摊），峰值按 $2$--$5$ 倍估。
单条记录按 $1$\,KB 估；$1$ 亿条 $\approx 100$\,GB。
内存读 $\sim$ 微秒级，SSD $\sim 0.1$\,ms，磁盘 $\sim 10$\,ms，跨地域网络 $\sim 100$\,ms。
\end{keypoint}

\section{一致性、CAP 与权衡}

分布式系统设计的取舍，绝大多数都能归结到“一致性 vs 可用性 vs 延迟”这条轴上。

\begin{itemize}
  \item \textbf{CAP}：网络分区（P）不可避免，因此真正的取舍是分区时选 \textbf{C}（一致性）还是 \textbf{A}（可用性）。订单/支付偏 C；信息流/搜索结果偏 A。
  \item \textbf{强一致 vs 最终一致}：强一致代价是延迟和可用性。\emph{什么时候必须强一致？}——涉及钱、涉及“不能重复占用的资源”（如同一房源同一晚）。其余大多数场景，最终一致足矣。
\end{itemize}

\begin{idea}[一致性是“分区域”选择的，不是全局的]
高手的设计不会笼统地说“整个系统强一致”。而是\textbf{按数据域分别选择}：预订的“占位”这一步必须强一致（防超卖）；而房源的浏览量、评分均值、搜索索引，最终一致即可。把强一致的范围\textbf{收缩到最小的关键路径}，是兼顾正确性与性能的核心手艺。
\end{idea}

\begin{followup}[“如果这里数据不一致了会怎样？”]
面试官常用这个问题探你对一致性的理解深度。应答框架：先说明\emph{这个数据域选了哪种一致性及原因}，再说\emph{不一致的具体后果}（用户看到旧评分 vs 一个房源被两人订走），最后给\emph{兜底手段}（对账任务、补偿事务、唯一约束作为最后防线）。
\end{followup}

\section{画架构图的小建议}

白板/共享文档上，用清晰的方框 + 箭头即可。建议固定从左到右的层次：\code{Client -> CDN/Gateway -> Service -> Cache -> DB}，旁路挂 \code{MQ}、\code{Search}、\code{Object Store}。后面几章的架构图都用 ASCII 框图给出，你可以照着练手绘。
